\documentclass[11pt,a4paper]{article}

% ============================================================
% Layout and typography
% ============================================================

\usepackage[
  a4paper,
  top=2.1cm,
  bottom=2.2cm,
  left=2.25cm,
  right=2.25cm
]{geometry}

\usepackage[T1]{fontenc}
\usepackage[utf8]{inputenc}
\usepackage{lmodern}
\usepackage{microtype}

\setlength{\parindent}{0pt}
\setlength{\parskip}{5pt}
\setlength{\emergencystretch}{2em}

% ============================================================
% Mathematics, tables, and lists
% ============================================================

\usepackage{booktabs}
\usepackage{tabularx}
\usepackage{array}
\usepackage{amsmath,amssymb}
\usepackage{enumitem}

\setlist[itemize]{
  leftmargin=1.25em,
  itemsep=1pt,
  topsep=3pt
}

\setlist[enumerate]{
  leftmargin=1.45em,
  itemsep=2pt,
  topsep=3pt
}

\newcolumntype{Y}{>{\raggedright\arraybackslash}X}

% ============================================================
% Colors and hyperlinks
% ============================================================

\usepackage{xcolor}
\usepackage{xurl}
\usepackage{hyperref}

\definecolor{navy}{HTML}{17365D}
\definecolor{darkgray}{HTML}{444444}
\definecolor{lightgray}{HTML}{F2F4F7}

\hypersetup{
  colorlinks=true,
  linkcolor=navy,
  citecolor=navy,
  urlcolor=navy,
  pdftitle={Data Construction Guide},
  pdfauthor={Momen Ismail},
  pdfsubject={Monthly stock return prediction data pipeline}
}

% ============================================================
% Header and footer
% ============================================================

\usepackage{fancyhdr}
\usepackage{lastpage}

\setlength{\headheight}{14pt}

\pagestyle{fancy}
\fancyhf{}
\lhead{Monthly Stock Return Prediction Project}
\rhead{Data Construction Guide}
\cfoot{\thepage\ of \pageref{LastPage}}
\renewcommand{\headrulewidth}{0.4pt}

% ============================================================
% Commands
% ============================================================

\newcommand{\code}[1]{\texttt{#1}}
\newcommand{\target}{\code{target\_excess\_return\_next\_1m}}

% ============================================================
% Document
% ============================================================

\begin{document}

\hypersetup{pageanchor=false}

% ============================================================
% Title page
% ============================================================

\begin{titlepage}
\centering

\vspace*{2.1cm}

{\Huge\bfseries\color{navy}
Data Construction Guide
\par}

\vspace{0.45cm}

{\Large
Monthly Stock Return Prediction Project
\par}

\vspace{1.0cm}

{\large
Final reproducibility, timing, cleaning, and validation documentation
\par}

\vfill

{\large
Momen Ismail\par
}

\vspace{0.2cm}

{\large
Matriculation Number: 50307564\par
}

\vspace{0.2cm}

{\large
Institutional Email:\par
\href{mailto:s73misma@uni-bonn.de}
{\texttt{s73misma@uni-bonn.de}}
\par}

\vspace{0.8cm}

{\large
Machine Learning for Finance\par
University of Bonn\par
Instructor: Prof. Ivan Gufler
\par}

\vfill

{\large\bfseries
Public Replication Repository
\par}

\vspace{0.25cm}

\url{https://github.com/Momen-Ismail/stock-return-ml-project}

\vfill

{\small
Final version -- July 2026
\par}

\vspace{0.1cm}

{\small
Final model sample: February 1990--December 2025
\par}

\end{titlepage}

% ============================================================
% Table of contents
% ============================================================

\tableofcontents
\newpage

\pagenumbering{arabic}
\hypersetup{pageanchor=true}

% ============================================================
% How to use the guide
% ============================================================

\section*{How to Use This Guide}
\addcontentsline{toc}{section}{How to Use This Guide}

This document is the detailed technical companion to the course-project report.
The main report presents the research question, model comparison, empirical
results, critical reflection, and conclusion. This guide explains exactly how
the common model-ready dataset was constructed.

The guide is linked directly from the repository's main \code{README.md}. A
reader who only wants to inspect the final empirical results does not need to
read the complete document. The most relevant sections are:

\begin{itemize}
  \item \textbf{Final Verified Dataset}, for the dimensions and final predictor
        architecture;
  \item \textbf{Pipeline Overview}, for the correct order of the five
        data-construction scripts;
  \item \textbf{Critical Design Decisions}, for the main methodological choices;
  \item \textbf{Timing and Leakage Controls}, for the information-availability
        rules;
  \item \textbf{Reproduction Steps}, for rebuilding the final dataset;
  \item \textbf{Limitations}, for the remaining data restrictions.
\end{itemize}

The detailed sections for Files 01--05 should be consulted when checking a
specific variable, formula, cleaning rule, lag, merge, input, or output file.

This guide documents the construction of the common model dataset. Model
estimation, hyperparameter tuning, validation-based specification selection,
and final-test results are documented in the main report and in the model-output
folders of the repository.

% ============================================================
% 1. Purpose
% ============================================================

\section{Purpose and Construction Contract}

This document records the final construction of the monthly stock-return
prediction dataset. Its purpose is to make the data pipeline reproducible,
auditable, and explicit about the most important methodological decisions.

Each observation is one stock-month. Predictors are dated at month \(t\), and
the target is the stock's excess return in month \(t+1\):

\[
  y_{i,t+1}=r_{i,t+1}-RF_{t+1}.
\]

The construction contract is:

\begin{itemize}
  \item one row per ticker-month;
  \item no use of information that becomes available after the predictor month;
  \item a raw and unwinsorized prediction target;
  \item monthly cross-sectional treatment of predictor missingness and outliers;
  \item no model-specific standardization or sample splitting inside the common
        data file;
  \item a frozen predictor list saved next to the final dataset.
\end{itemize}

The final model-ready dataset is:

\begin{center}
\path{output/data/final/006_model_dataset_kelly_winsorized_1990_2025.parquet}
\end{center}

The authoritative predictor list is:

\begin{center}
\path{output/data/final/006_predictor_columns_kelly_winsorized.csv}
\end{center}

The final dataset contains:

\begin{itemize}
  \item \code{ticker};
  \item \code{month};
  \item the raw target \target;
  \item 484 predictors.
\end{itemize}

Chronological training, validation, development, and test samples are created
later in the modeling workflow rather than stored as separate data-construction
files.

% ============================================================
% 2. Final dataset
% ============================================================

\section{Final Verified Dataset}

After correcting the Welch--Goyal source and rerunning Files 03--05, the final
verified dataset has:

\begin{center}
\begin{tabularx}{0.86\textwidth}{p{6.0cm}Y}
\toprule
Item & Verified value \\
\midrule
Rows & 211,059 \\
Total columns & 487 \\
Predictors & 484 \\
Unique tickers & 656 \\
First model month & 1990-02-28 \\
Last model month & 2025-12-31 \\
Missing predictor values & 0 \\
Infinite predictor values & 0 \\
Duplicate ticker-month rows & 0 \\
Base predictors & 124 \\
Continuous stock characteristics & 45 \\
Binary stock characteristics & 3 \\
Market and VIX variables & 4 \\
Welch--Goyal macro variables & 8 \\
SIC2 dummies & 64 \\
Characteristic--macro interactions & 360 \\
\bottomrule
\end{tabularx}
\end{center}

The 124 base predictors are composed of:

\[
45
+3
+4
+8
+64
=124.
\]

More specifically, they contain:

\begin{itemize}
  \item 45 continuous stock characteristics;
  \item 3 binary stock characteristics;
  \item 4 market/VIX variables;
  \item 8 macroeconomic states;
  \item 64 SIC2 industry indicators.
\end{itemize}

The 360 interactions are the Cartesian product of the 45 continuous stock
characteristics and the 8 macroeconomic states:

\[
45\times 8=360.
\]

Binary characteristics, market variables, macro variables, and industry dummies
are not interacted.

% ============================================================
% 3. Pipeline overview
% ============================================================

\section{Pipeline Overview}

\begin{center}
\small
\begin{tabularx}{\textwidth}{p{1.7cm}p{5.3cm}Y}
\toprule
Stage & Main script & Role \\
\midrule

File 01
&
\path{src/data/01_build_clean_yahoo_daily.py}
&
Builds and cleans the daily Yahoo stock-price panel for the locked stock
universe.
\\

File 02
&
\path{src/data/02_build_monthly_stock_features.py}
&
Aggregates daily observations to monthly predictors, adds market/VIX
information, and constructs the next-month excess-return target.
\\

File 03
&
\path{src/data/03_add_fundamentals_and_macro.py}
&
Cleans annual Compustat data and merges only reports considered available after
a conservative six-month lag. Despite the historical filename, Welch--Goyal
variables are not constructed in this file.
\\

File 04
&
\path{src/data/04_build_raw_kelly_dataset.py}
&
Adds SIC2 dummies, merges one-month-lagged Welch--Goyal variables, validates the
macroeconomic source, and freezes the 124 base predictors.
\\

File 05
&
\path{src/data/05_clean_and_rank_normalize.py}
&
Performs monthly median imputation and monthly 1st/99th percentile
winsorization, then constructs the 360 interactions. The filename is retained
for compatibility, but the final data are not rank-normalized.
\\

\bottomrule
\end{tabularx}
\end{center}

The five files must be run in numerical order. Each stage reads the output of
the previous stage and saves a new auditable file.

% ============================================================
% 4. Inputs
% ============================================================

\section{Permanent and Restricted Inputs}

\subsection{Version-controlled reproducibility inputs}

The following inputs are intended to remain under version control:

\begin{center}
\small
\begin{tabularx}{\textwidth}{p{6.0cm}Y}
\toprule
File & Role \\
\midrule

\path{input/stock_universe_locked.csv}
&
Locked Yahoo-format stock universe used by File 01.
\\

\path{input/market_gspc_daily.csv}
&
Permanent daily S\&P 500 market series used for market returns, market
volatility, beta, and idiosyncratic volatility.
\\

\path{input/market_vix_daily.csv}
&
Permanent daily VIX series used to construct monthly VIX predictors.
\\

\path{input/external/fama_french_rf_monthly.csv}
&
Monthly risk-free rate used to construct the excess-return target.
\\

\path{input/external/welch_goyal_macro_1990_2025.csv}
&
Cleaned monthly macro predictors used by File 04.
\\

\path{input/raw/PredictorData2025.xlsx}
&
Original updated Welch--Goyal workbook used to reproduce the cleaned macro CSV.
\\

\path{input/input_manifest.csv}
&
Input registry with coverage and source notes.
\\

\path{input/README.md}
&
Input-folder documentation.
\\

\bottomrule
\end{tabularx}
\end{center}

The obsolete file:

\begin{center}
\path{input/raw/PredictorData.xls}
\end{center}

ended in December 2005 and is not used by the final project.

\subsection{Restricted or externally acquired inputs}

Annual Compustat data are licensed and are not intended for public
redistribution. The raw Compustat location is defined in:

\begin{center}
\path{src/config.py}
\end{center}

A complete reconstruction therefore requires the same WRDS/Compustat extraction
with the variables documented in Section~\ref{sec:compustat}.

The public repository contains the construction code, variable definitions,
cleaning rules, six-month availability lag, and final output structure, even
when the raw licensed file cannot be distributed.

Daily Yahoo stock prices are acquired or loaded by File 01 for the locked
ticker universe. The locked universe and quality reports provide a reproducible
record of the intended sample. However, Yahoo symbols are not permanent
security identifiers, and historical vendor availability may change.

% ============================================================
% 5. Design decisions
% ============================================================

\section{Critical Design Decisions}

\subsection{Locked stock universe}

The project uses a locked S\&P 500-style ticker universe rather than
redownloading a changing list during each run.

This decision prevents the list of firms from changing silently across
executions. It improves internal consistency and reproducibility.

However, it does not eliminate survivorship or index-selection concerns because
the project does not use a fully point-in-time CRSP universe with permanent
security identifiers and historical constituent membership.

\subsection{Time-aware target and predictors}

Predictors in month \(t\) are used to forecast excess return in month \(t+1\).
Chronological ordering is enforced before all group-wise shifts and rolling
calculations.

The target is:

\[
  \target_{i,t}
  =
  r_{i,t+1}-RF_{t+1}.
\]

The target is retained only when the next observation for the ticker is the
exact next calendar month. A return spanning several calendar months is not
treated as a one-month return.

The target is not winsorized in the common dataset. Large realized returns are
genuine outcomes and are retained for transparent evaluation.

Any model-specific target treatment would have to be estimated only on the
relevant training sample and reported separately. It is not part of the
official results.

\subsection{No global preprocessing leakage}

The common data file does not use full-sample standardization.

Models that require scaling, including OLS-3, Elastic Net, and PLS, standardize
predictors later inside training-only model pipelines.

Tree-based models use the common predictors without standardization.

Chronological validation is performed with expanding annual folds rather than
random cross-validation.

\subsection{Monthly predictor treatment}

File 05 applies the following predictor treatment:

\begin{enumerate}
  \item continuous stock characteristics are imputed using the same-month
        cross-sectional median;
  \item continuous stock characteristics are winsorized at the same-month 1st
        and 99th percentiles;
  \item binary characteristics are filled with zero when missing and remain
        binary;
  \item market variables, macro variables, and SIC2 dummies are not
        rank-normalized;
  \item residual missing values after the defined treatment are filled with
        zero;
  \item the target is left unchanged.
\end{enumerate}

The resulting predictors remain in economically interpretable units, subject
only to monthly winsorization of the continuous stock-characteristic block.

\subsection{Interaction design}

The project uses a compact Kelly-style conditional-predictability design.

Only continuous stock characteristics are interacted with the eight
Welch--Goyal macro states:

\[
  Interaction_{i,t}^{c,m}
  =
  Characteristic_{i,t}^{c}
  \times
  MacroState_t^{m}.
\]

The project intentionally excludes:

\begin{itemize}
  \item binary-characteristic \(\times\) macro interactions;
  \item market/VIX \(\times\) macro interactions;
  \item macro \(\times\) macro interactions;
  \item SIC2 \(\times\) macro interactions;
  \item interactions involving the target.
\end{itemize}

This decision keeps the feature space manageable and interpretable while still
allowing characteristic effects to vary with aggregate market conditions.

% ============================================================
% 6. File 01
% ============================================================

\section{File 01: Daily Stock Prices and Quality Cleaning}

\textbf{Main script:}

\begin{center}
\path{src/data/01_build_clean_yahoo_daily.py}
\end{center}

\textbf{Purpose:} create a clean daily OHLCV stock panel for the locked stock
universe.

File 01 downloads or reads Yahoo data with \code{auto\_adjust=False},
standardizes ticker and field names, sorts observations chronologically, and
computes daily adjusted returns:

\[
  r_{i,d}
  =
  \frac{AdjClose_{i,d}}
       {AdjClose_{i,d-1}}
  -1.
\]

Daily quality rules identify invalid observations involving:

\begin{itemize}
  \item missing core OHLCV fields;
  \item nonpositive adjusted prices;
  \item impossible OHLC relations, such as \code{high < low};
  \item duplicate ticker-date rows;
  \item extreme daily returns above 300\% or below \(-95\%\).
\end{itemize}

A complete ticker history is removed only when its bad-row share exceeds the
configured threshold. This avoids deleting an otherwise usable history because
of a small number of isolated problematic observations.

File 01 saves:

\begin{itemize}
  \item the cleaned daily stock-price panel;
  \item the final clean ticker list;
  \item removed-ticker information;
  \item data-quality reports.
\end{itemize}

% ============================================================
% 7. File 02
% ============================================================

\section{File 02: Monthly Stock Features and Target}

\textbf{Main script:}

\begin{center}
\path{src/data/02_build_monthly_stock_features.py}
\end{center}

\textbf{Purpose:} aggregate daily stock, market, VIX, and risk-free-rate inputs
into a monthly stock panel.

Monthly return is calculated from month-end adjusted prices:

\[
  r_{i,t}
  =
  \frac{AdjClose_{i,t}}
       {AdjClose_{i,t-1}}
  -1.
\]

Momentum variables use only lagged monthly returns and exclude the
contemporaneous month:

\[
  mom_{i,t}^{(k)}
  =
  \prod_{j=1}^{k}
  (1+r_{i,t-j})-1.
\]

The retained monthly stock predictors include:

\begin{itemize}
  \item recent return;
  \item six- and twelve-month momentum;
  \item long-horizon return;
  \item realized volatility;
  \item downside-return and maximum-return measures;
  \item dollar volume;
  \item Amihud illiquidity;
  \item zero-trading activity;
  \item price-trend measures;
  \item beta;
  \item idiosyncratic volatility.
\end{itemize}

Amihud illiquidity is calculated from daily data as:

\[
  Amihud_{i,d}
  =
  \frac{|r_{i,d}|}
  {AdjClose_{i,d}\times Volume_{i,d}},
\]

and is then averaged within each month.

\subsection{Zero-volume correction}

A final audit identified 136 daily observations with invalid zero volume.

The correction was deliberately narrow. The observations were not deleted and
unrelated price predictors were not altered.

Only volume-derived monthly quantities were set to missing when the underlying
daily volume information was invalid:

\begin{itemize}
  \item \code{avg\_dolvol\_1m};
  \item \code{std\_dolvol\_1m};
  \item \code{avg\_log\_dolvol\_1m};
  \item \code{amihud\_1m};
  \item \code{zerotrade\_1m}.
\end{itemize}

The variable \code{dolvol\_growth\_1m} becomes missing automatically when its
underlying monthly dollar-volume input is missing.

These missing values are later handled by the same-month median-imputation rule
in File 05.

This correction preserves valid return and price information while preventing
an invalid volume observation from contaminating liquidity predictors.

\subsection{Market, VIX, beta, and idiosyncratic volatility}

The direct market-state predictors are:

\[
\code{market\_ret\_1m},
\quad
\code{market\_vol\_1m},
\quad
\code{vix\_avg\_1m},
\quad
\code{vix\_change\_1m}.
\]

Rolling beta and idiosyncratic volatility use daily stock and market returns
over a 252-trading-day window with a minimum of 126 observations:

\[
  \beta_{i,t}
  =
  \frac{
    \operatorname{Cov}(r_{i,d},r_{m,d})
  }{
    \operatorname{Var}(r_{m,d})
  },
\]

\[
  \sigma^{idio}_{i,t}
  =
  \sqrt{
    \max
    \left(
      \operatorname{Var}(r_{i,d})
      -
      \beta_{i,t}^{2}
      \operatorname{Var}(r_{m,d}),
      0
    \right)
  }.
\]

The locked risk block also includes:

\[
\code{betasq\_12m}
=
\code{beta\_12m}^{2}.
\]

\subsection{Target construction}

The next-month target is calculated within ticker after chronological sorting.

For a predictor row in month \(t\), the target is the excess return realized in
month \(t+1\):

\[
  y_{i,t+1}
  =
  r_{i,t+1}
  -
  RF_{t+1}.
\]

The row is retained only when the next ticker observation is the exact next
calendar month.

This rule prevents an observation after a long gap from being mislabeled as a
one-month return.

% ============================================================
% 8. File 03
% ============================================================

\section{File 03: Annual Compustat Characteristics and Timing}
\label{sec:compustat}

\textbf{Main script:}

\begin{center}
\path{src/data/03_add_fundamentals_and_macro.py}
\end{center}

\textbf{Input:}

\begin{itemize}
  \item monthly stock panel;
  \item raw annual Compustat file.
\end{itemize}

\textbf{Outputs:}

\begin{itemize}
  \item \path{output/data/staging/003_compustat_annual_clean.parquet};
  \item \path{output/data/staging/004_monthly_panel_with_fundamentals.parquet}.
\end{itemize}

Compustat \code{datadate} is the fiscal-year-end date, not the date on which the
report became publicly available.

The project therefore imposes a conservative six-month availability lag:

\[
  AvailableMonth
  =
  MonthEnd(datadate+6\text{ months}).
\]

For each ticker-month, a backward as-of merge selects the most recent Compustat
observation whose availability month is not later than the stock month.

Reports older than \code{MAX\_COMPUSTAT\_AGE\_MONTHS} are invalidated as stale.
This prevents an old report from being carried forward indefinitely.

The latest verified File 03 results are:

\begin{itemize}
  \item clean annual Compustat rows: 21,575;
  \item clean Compustat tickers: 649;
  \item monthly panel rows: 211,341;
  \item monthly panel tickers: 656;
  \item matched observations: 207,213;
  \item unmatched observations: 4,128;
  \item stale matches invalidated: 353;
  \item timing violations: 0;
  \item missing targets: 0.
\end{itemize}

\subsection{Accounting variables}

\begin{center}
\small
\begin{tabularx}{\textwidth}{p{4.0cm}Y}
\toprule
Family & Retained variables \\
\midrule

Valuation and size
&
\code{be\_me},
\code{ocf\_me},
\code{log\_comp\_market\_equity}
\\

Profitability
&
\code{op\_at},
\code{ocf\_at}
\\

Leverage and liquidity
&
\code{debt\_at},
\code{cash\_at},
\code{cashflow\_to\_debt},
\code{current\_ratio},
\code{quick\_ratio}
\\

Investment and asset structure
&
\code{capx\_at},
\code{capx\_growth},
\code{rd\_at},
\code{ppe\_at},
\code{tangibility},
\code{asset\_turnover},
\code{sales\_to\_inventory},
\code{sales\_to\_cash},
\code{sales\_to\_receivables}
\\

Accruals, payout, and financing
&
\code{accruals\_at},
\code{dividend\_dummy},
\code{equity\_issuance\_at},
\code{repurchase\_at}
\\

Growth and missingness
&
\code{asset\_growth},
\code{sales\_growth},
\code{ppeinv\_gr1a},
\code{xrd\_missing},
\code{has\_compustat\_annual}
\\

Industry
&
\code{sic2} before dummy encoding
\\

\bottomrule
\end{tabularx}
\end{center}

Annual growth rates are calculated only when consecutive reports are
approximately one year apart.

Gaps outside 300--430 days cause the affected growth variables to be set to
missing rather than treating irregularly spaced reports as one-year growth.

% ============================================================
% 9. Welch-Goyal
% ============================================================

\section{Welch--Goyal Input Reconstruction and Validation}

A final variable-range audit revealed that the earlier cleaned macro file
contained only 124 unique values for both \code{wg\_dp} and \code{wg\_ep}
across 431 model months.

Long exact constant sequences were present, including a 197-month block from
August 2009 through December 2025.

The merge in File 04 was correct. The problem originated in the previously
prepared external input.

The obsolete source workbook:

\begin{center}
\path{input/raw/PredictorData.xls}
\end{center}

ended in December 2005.

The corrected reproducibility source is:

\begin{center}
\path{input/raw/PredictorData2025.xlsx}
\end{center}

Its monthly sheet contains 1,860 observations from January 1871 through
December 2025 and includes:

\begin{itemize}
  \item \code{Index};
  \item \code{D12};
  \item \code{E12};
  \item \code{b/m};
  \item \code{tbl};
  \item \code{AAA};
  \item \code{BAA};
  \item \code{lty};
  \item \code{ntis};
  \item \code{svar}.
\end{itemize}

The acquisition script:

\begin{center}
\path{src/acquisition/04_create_welch_goyal_input.py}
\end{center}

reconstructs the cleaned macro file for January 1990--December 2025 using:

\[
  wg\_dp_t
  =
  \log
  \left(
  \frac{D12_t}{Index_t}
  \right),
\]

\[
  wg\_ep_t
  =
  \log
  \left(
  \frac{E12_t}{Index_t}
  \right),
\]

\[
  wg\_tms_t
  =
  lty_t-tbl_t,
\]

\[
  wg\_dfy_t
  =
  BAA_t-AAA_t.
\]

The remaining variables are copied directly:

\[
  wg\_bm=b/m,
  \qquad
  wg\_ntis=ntis,
\]

\[
  wg\_tbl=tbl,
  \qquad
  wg\_svar=svar.
\]

No forward filling is used.

The validator checks:

\begin{itemize}
  \item valid and unique monthly dates;
  \item a continuous monthly sequence;
  \item numeric macro columns without missing or infinite values;
  \item suspicious exact constant runs for \code{wg\_dp} and \code{wg\_ep}.
\end{itemize}

The corrected external file has 432 monthly observations and 432 unique values
for both \code{wg\_dp} and \code{wg\_ep}.

After applying the one-month lag and removing January 1990, the final model
panel contains 431 unique monthly values for each variable.

% ============================================================
% 10. File 04
% ============================================================

\section{File 04: Raw Kelly-Style Base Panel}

\textbf{Main script:}

\begin{center}
\path{src/data/04_build_raw_kelly_dataset.py}
\end{center}

\textbf{Inputs:}

\begin{itemize}
  \item \path{output/data/staging/004_monthly_panel_with_fundamentals.parquet};
  \item \path{input/external/welch_goyal_macro_1990_2025.csv}.
\end{itemize}

\textbf{Outputs:}

\begin{itemize}
  \item \path{output/data/staging/005_raw_kelly_base.parquet};
  \item \path{output/data/staging/005_raw_kelly_predictors.csv};
  \item \path{output/quality/005_raw_kelly_summary.csv}.
\end{itemize}

File 04 creates SIC2 industry dummies, including
\code{sic2\_missing}, and requires exactly one industry category per
observation.

It then shifts each Welch--Goyal source month forward by one month before
merging.

Therefore, macro information dated month \(t\) becomes a predictor for stock
month \(t+1\).

January 1990 has no December 1989 macro source row in the restricted cleaned
file. Its 282 stock observations therefore contain missing macro states.

This is expected. File 05 removes this complete first month, yielding a final
sample beginning in February 1990.

The verified File 04 base panel contains:

\begin{itemize}
  \item rows: 211,341;
  \item total columns: 127;
  \item tickers: 656;
  \item base predictors: 124;
  \item continuous stock characteristics: 45;
  \item binary stock characteristics: 3;
  \item market variables: 4;
  \item macro variables: 8;
  \item SIC2 dummies: 64;
  \item rows with missing macro states: 282, all in January 1990;
  \item missing targets: 0;
  \item infinite values: 0;
  \item duplicate ticker-month rows: 0.
\end{itemize}

% ============================================================
% 11. File 05
% ============================================================

\section{File 05: Imputation, Winsorization, and Interactions}

\textbf{Main script:}

\begin{center}
\path{src/data/05_clean_and_rank_normalize.py}
\end{center}

\textbf{Inputs:}

\begin{itemize}
  \item \path{output/data/staging/005_raw_kelly_base.parquet};
  \item \path{output/data/staging/005_raw_kelly_predictors.csv}.
\end{itemize}

\textbf{Outputs:}

\begin{itemize}
  \item \path{output/data/final/006_model_dataset_kelly_winsorized_1990_2025.parquet};
  \item \path{output/data/final/006_predictor_columns_kelly_winsorized.csv};
  \item \path{output/quality/006_cleaning_summary.csv}.
\end{itemize}

The script name still contains \code{rank\_normalize} for compatibility with the
earlier project structure.

The final methodology does not rank-normalize the predictors.

Instead, File 05 applies:

\begin{itemize}
  \item monthly cross-sectional median imputation;
  \item monthly 1st/99th percentile winsorization of continuous stock
        characteristics;
  \item binary-value preservation;
  \item interaction construction;
  \item final data-quality checks.
\end{itemize}

\subsection{Variable blocks}

The 48 stock characteristics are divided into:

\begin{itemize}
  \item 45 continuous characteristics;
  \item 3 binary characteristics:
        \code{dividend\_dummy},
        \code{xrd\_missing}, and
        \code{has\_compustat\_annual}.
\end{itemize}

Continuous characteristics are imputed and winsorized within month.

Binary variables are filled with zero when missing and remain in
\(\{0,1\}\).

The following blocks are preserved without cross-sectional winsorization:

\begin{itemize}
  \item 4 market/VIX variables;
  \item 8 macro variables;
  \item 64 SIC2 dummies.
\end{itemize}

After characteristic processing, File 05 constructs 360 continuous
characteristic \(\times\) macro interactions.

The final file contains 484 predictors and 487 total columns.

% ============================================================
% 12. Predictor groups
% ============================================================

\section{Final Predictor Groups}

\begin{center}
\small
\begin{tabularx}{\textwidth}{p{3.8cm}Y}
\toprule
Group & Examples and retained variables \\
\midrule

Recent return and momentum
&
\code{ret\_1m},
\code{mom6m},
\code{mom12m},
\code{long\_term\_return\_36m}
\\

Volatility and tail risk
&
\code{retvol\_1m},
\code{retvol12m},
\code{minret\_1m},
\code{rmax5\_21d},
\code{beta\_12m},
\code{betasq\_12m},
\code{idiovol\_12m}
\\

Liquidity and trading activity
&
\code{avg\_dolvol\_1m},
\code{std\_dolvol\_1m},
\code{avg\_log\_dolvol\_1m},
\code{amihud\_1m},
\code{zerotrade\_1m},
\code{dolvol\_growth\_1m}
\\

Price trend
&
\code{price\_to\_ma12},
\code{dist\_from\_high\_12m},
\code{avg\_range\_1m}
\\

Market and VIX state
&
\code{market\_ret\_1m},
\code{market\_vol\_1m},
\code{vix\_avg\_1m},
\code{vix\_change\_1m}
\\

Valuation and size
&
\code{be\_me},
\code{ocf\_me},
\code{log\_comp\_market\_equity}
\\

Profitability
&
\code{op\_at},
\code{ocf\_at}
\\

Leverage and liquidity
&
\code{debt\_at},
\code{cash\_at},
\code{cashflow\_to\_debt},
\code{current\_ratio},
\code{quick\_ratio}
\\

Investment and asset structure
&
\code{capx\_at},
\code{capx\_growth},
\code{rd\_at},
\code{ppe\_at},
\code{tangibility},
\code{asset\_turnover},
\code{sales\_to\_inventory},
\code{sales\_to\_cash},
\code{sales\_to\_receivables}
\\

Accruals, payout, and financing
&
\code{accruals\_at},
\code{dividend\_dummy},
\code{equity\_issuance\_at},
\code{repurchase\_at}
\\

Growth and missingness
&
\code{asset\_growth},
\code{sales\_growth},
\code{ppeinv\_gr1a},
\code{xrd\_missing},
\code{has\_compustat\_annual}
\\

Industry controls
&
64 SIC2 dummies, including \code{sic2\_missing}
\\

Macro states
&
\code{wg\_dp},
\code{wg\_ep},
\code{wg\_bm},
\code{wg\_ntis},
\code{wg\_tbl},
\code{wg\_tms},
\code{wg\_dfy},
\code{wg\_svar}
\\

Interactions
&
Each of the 45 continuous stock characteristics multiplied by each of the
8 macro states
\\

\bottomrule
\end{tabularx}
\end{center}

% ============================================================
% 13. Timing controls
% ============================================================

\section{Timing and Leakage Controls}

\begin{center}
\small
\begin{tabularx}{\textwidth}{p{4.0cm}Y}
\toprule
Control & Implementation \\
\midrule

Chronological sorting
&
Data are sorted by ticker and date before all shifts and rolling calculations.
\\

Target timing
&
A month-\(t\) row predicts the stock's excess return in month \(t+1\).
\\

Consecutive-month condition
&
The target is retained only when the next ticker observation is the exact next
calendar month.
\\

Momentum timing
&
Momentum variables use lagged monthly returns and exclude the contemporaneous
month.
\\

Market timing
&
Monthly market and VIX predictors use information observed by the end of month
\(t\).
\\

Macro timing
&
Welch--Goyal month \(t\) is shifted forward and used for stock month \(t+1\).
\\

Accounting timing
&
Annual Compustat information becomes available six months after fiscal-year
end.
\\

Accounting merge
&
A backward as-of merge prevents future reports from entering past stock months.
\\

Stale reports
&
Accounting matches older than the configured maximum age are invalidated.
\\

Imputation
&
Continuous characteristics use only the same month's cross-section.
\\

Winsorization
&
Monthly 1st/99th percentiles are calculated only within the contemporaneous
cross-section.
\\

Model scaling
&
Standardization is deferred to training-only model pipelines.
\\

Model validation
&
Annual expanding-window folds are used rather than random cross-validation.
\\

Final test
&
The 2020--2025 test period is held out from tuning and model selection.
\\

\bottomrule
\end{tabularx}
\end{center}

% ============================================================
% 14. Validation
% ============================================================

\section{Validation Results}

The final mechanical validation confirms:

\begin{itemize}
  \item zero missing predictor values;
  \item zero infinite predictor values;
  \item zero duplicate ticker-month rows;
  \item zero constant predictors;
  \item unique predictor names;
  \item each SIC2 row sums to exactly one industry dummy;
  \item each market and macro variable has one value per month;
  \item all 360 interaction columns equal the exact product of their processed
        characteristic and macro components;
  \item \code{wg\_dp} and \code{wg\_ep} have 431 distinct values across the
        431 final model months;
  \item the December 2025 stock rows contain November 2025 Welch--Goyal values,
        confirming the intended one-month lag.
\end{itemize}

The raw target is intentionally not clipped.

The variable-range audit found a minimum of approximately \(-0.935\) and a
maximum of approximately \(16.251\).

These observations are retained because the target represents realized
next-month excess returns.

Evaluation therefore reports MAE as a robust complement to squared-error
measures, although monthly MSE remains the main project metric.

% ============================================================
% 15. Reproduction
% ============================================================

\section{Reproduction Steps}

All commands in this section must be run from the repository root.

The repository root is the folder containing:

\begin{itemize}
  \item \code{README.md};
  \item \code{requirements.txt};
  \item \code{src/};
  \item \code{input/};
  \item \code{output/};
  \item \code{documentation/}.
\end{itemize}

A complete reconstruction requires the licensed Compustat input. Because this
file cannot be redistributed publicly, it must be placed at the location
configured in:

\begin{center}
\path{src/config.py}
\end{center}

before running File 03.

The numbered commands below are the authoritative execution order.

The five construction scripts should be run directly rather than relying on an
obsolete pipeline wrapper.

\subsection{Install Python requirements}

From the repository root:

\begin{center}
\code{python3 -m pip install -r requirements.txt}
\end{center}

A separate virtual environment is recommended.

\subsection{Rebuild the Welch--Goyal input when necessary}

If the cleaned Welch--Goyal CSV needs to be reconstructed from the permanent
Excel source, run:

\begin{center}
\code{python3 src/acquisition/04\_create\_welch\_goyal\_input.py}
\end{center}

This script is deterministic given:

\begin{center}
\path{input/raw/PredictorData2025.xlsx}
\end{center}

\subsection{Run the five data-construction stages}

Run:

\begin{enumerate}
  \item
  \code{python3 src/data/01\_build\_clean\_yahoo\_daily.py}

  \item
  \code{python3 src/data/02\_build\_monthly\_stock\_features.py}

  \item
  \code{python3 src/data/03\_add\_fundamentals\_and\_macro.py}

  \item
  \code{python3 src/data/04\_build\_raw\_kelly\_dataset.py}

  \item
  \code{python3 src/data/05\_clean\_and\_rank\_normalize.py}
\end{enumerate}

The normal production scripts should not silently replace permanent market,
VIX, Fama--French, or Welch--Goyal inputs.

\subsection{Verify the final outputs}

After a successful run, the main final files should exist at:

\begin{center}
\path{output/data/final/006_model_dataset_kelly_winsorized_1990_2025.parquet}
\end{center}

\begin{center}
\path{output/data/final/006_predictor_columns_kelly_winsorized.csv}
\end{center}

\begin{center}
\path{output/quality/006_cleaning_summary.csv}
\end{center}

The final quality summary should report:

\begin{itemize}
  \item 211,059 observations;
  \item 656 tickers;
  \item 484 predictors;
  \item zero missing predictor values;
  \item zero infinite predictor values;
  \item zero duplicate ticker-month rows.
\end{itemize}

% ============================================================
% 16. Modeling interface
% ============================================================

\section{Modeling Interface}

The predictor file is authoritative.

All modeling scripts load predictors from:

\begin{center}
\path{output/data/final/006_predictor_columns_kelly_winsorized.csv}
\end{center}

The predictor list should not be reconstructed manually inside individual
model scripts.

This ensures that:

\begin{itemize}
  \item every model receives the same 484 predictors;
  \item obsolete variables cannot enter the comparison;
  \item intermediate construction columns cannot be selected accidentally;
  \item the final model inputs remain reproducible.
\end{itemize}

\subsection{Chronological samples}

The final sample design is:

\begin{center}
\small
\begin{tabularx}{0.94\textwidth}{p{3.2cm}Y}
\toprule
Sample & Period and role \\
\midrule

Training
&
February 1990--December 2014. Used for fixed-model estimation and
hyperparameter tuning.
\\

Validation
&
January 2015--December 2019. Used to compare fixed and optimized
specifications.
\\

Development
&
February 1990--December 2019. Used to refit the locked final model
specifications.
\\

Test
&
January 2020--December 2025. Used once for official final evaluation.
\\

\bottomrule
\end{tabularx}
\end{center}

The corresponding final sample sizes are:

\begin{center}
\small
\begin{tabularx}{0.94\textwidth}{p{2.7cm}r r r Y}
\toprule
Sample & Observations & Months & Tickers & Period \\
\midrule
Training & 131,061 & 299 & 572 & Feb. 1990--Dec. 2014 \\
Validation & 35,386 & 60 & 611 & Jan. 2015--Dec. 2019 \\
Development & 166,447 & 359 & 613 & Feb. 1990--Dec. 2019 \\
Test & 44,612 & 72 & 641 & Jan. 2020--Dec. 2025 \\
\bottomrule
\end{tabularx}
\end{center}

Ticker counts differ across periods because firms enter and leave the available
panel over time.

\subsection{Expanding-window tuning}

Hyperparameter tuning uses ten annual expanding-window folds inside the
training period:

\begin{center}
\small
\begin{tabularx}{0.82\textwidth}{p{4.6cm}Y}
\toprule
Training period & Validation year \\
\midrule
Through December 2004 & 2005 \\
Through December 2005 & 2006 \\
Through December 2006 & 2007 \\
Through December 2007 & 2008 \\
Through December 2008 & 2009 \\
Through December 2009 & 2010 \\
Through December 2010 & 2011 \\
Through December 2011 & 2012 \\
Through December 2012 & 2013 \\
Through December 2013 & 2014 \\
\bottomrule
\end{tabularx}
\end{center}

The primary tuning criterion is average fold-level monthly mean squared error.

Random cross-validation is not used.

\subsection{Final model set}

The final active model set is:

\begin{itemize}
  \item Historical Mean;
  \item OLS-3;
  \item Partial Least Squares;
  \item Elastic Net;
  \item Random Forest;
  \item Gradient Boosting.
\end{itemize}

A Decision Tree is not part of the official final comparison.

\subsection{Specification-selection procedure}

The final specifications were selected using the following procedure:

\begin{enumerate}
  \item estimate fixed specifications on the training sample;
  \item tune model families using the 2005--2014 expanding folds;
  \item estimate the tuned candidates on the complete training sample;
  \item compare fixed and optimized versions on the separate 2015--2019
        validation period;
  \item lock the final specification for each model family;
  \item refit the locked models on the 1990--2019 development sample;
  \item evaluate once on the untouched 2020--2025 test period.
\end{enumerate}

The final validation-based choices are:

\begin{center}
\small
\begin{tabularx}{0.94\textwidth}{p{3.2cm}Y}
\toprule
Model family & Final choice \\
\midrule

Partial Least Squares
&
Optimized specification with 2 components.
\\

Elastic Net
&
Optimized specification with
\code{alpha = 0.015} and
\code{l1\_ratio = 0.85}.
\\

Random Forest
&
Fixed specification with 100 trees.
\\

Gradient Boosting
&
Fixed specification with 100 trees, learning rate 0.01, and maximum depth 2.
\\

\bottomrule
\end{tabularx}
\end{center}

The optimized 300-tree Random Forest and Gradient Boosting candidates were not
used in the final test because they did not improve validation monthly MSE
relative to the fixed 100-tree specifications.

\subsection{Model-specific scaling}

OLS-3, PLS, and Elastic Net use standardized predictors.

Standardization is implemented inside the model-estimation pipeline and is
fitted only on the relevant training sample.

Random Forest and Gradient Boosting do not require predictor standardization
because their splits depend on value ordering rather than measurement units.

\subsection{Test-sample rule}

The final test period must not be used for:

\begin{itemize}
  \item hyperparameter tuning;
  \item predictor selection;
  \item choosing between fixed and optimized specifications;
  \item replacing one model with another;
  \item changing a specification after observing test performance.
\end{itemize}

The final 2020--2025 test is used once after all model decisions have been
locked.

% ============================================================
% 17. Relation to README and paper
% ============================================================

\section{Relationship to the README and Main Report}

The three main project documents serve different purposes.

\subsection{README}

The repository's main:

\begin{center}
\path{README.md}
\end{center}

is the operational entry point.

It explains:

\begin{itemize}
  \item where the important files are located;
  \item how to clone and inspect the repository;
  \item how to install dependencies;
  \item how to run the data and model stages;
  \item where the final outputs are stored;
  \item how to rebuild the report.
\end{itemize}

\subsection{Data Construction Guide}

This document provides the detailed technical reference for:

\begin{itemize}
  \item input sources;
  \item daily and monthly filters;
  \item predictor definitions;
  \item accounting availability;
  \item macroeconomic timing;
  \item imputation;
  \item winsorization;
  \item interactions;
  \item validation checks;
  \item data limitations.
\end{itemize}

\subsection{Main report}

The main course-project report is stored at:

\begin{center}
\path{documentation/pdf/Momen Ismail.pdf}
\end{center}

Its editable source is:

\begin{center}
\path{documentation/documents/thesis/paper.tex}
\end{center}

The report presents:

\begin{itemize}
  \item the research question;
  \item the model-comparison design;
  \item the final empirical results;
  \item the interpretation;
  \item the critical reflection;
  \item the conclusion.
\end{itemize}

The report should be used for the project's empirical conclusions. This guide
should be used when checking how a particular observation, variable, lag,
cleaning decision, or data file was constructed.

% ============================================================
% 18. Build guide
% ============================================================

\section{Building This Guide}

The LaTeX source is stored at:

\begin{center}
\path{documentation/documents/data_construction/Data_Construction_Guide.tex}
\end{center}

To compile the guide directly, run:

\begin{verbatim}
cd documentation/documents/data_construction
latexmk -pdf -interaction=nonstopmode -halt-on-error \
Data_Construction_Guide.tex
mkdir -p ../../pdf
cp Data_Construction_Guide.pdf \
../../pdf/Data_Construction_Guide.pdf
cd ../../..
\end{verbatim}

The compiled PDF should then be available at:

\begin{center}
\path{documentation/pdf/Data_Construction_Guide.pdf}
\end{center}

The README should link directly to this PDF so that it can be opened from the
GitHub repository without compiling the LaTeX source.

% ============================================================
% 19. Limitations
% ============================================================

\section{Limitations}

The final dataset is designed for a transparent course project rather than a
complete replication of the CRSP/Compustat universe used in the largest
empirical asset-pricing studies.

The main limitations are:

\begin{itemize}
  \item the locked S\&P 500-style universe creates survivorship and
        index-selection concerns;
  \item Yahoo ticker histories are not permanent security identifiers;
  \item ticker changes, mergers, acquisitions, and delistings may complicate
        historical matching;
  \item Yahoo Finance does not provide complete delisting returns;
  \item exact Compustat filing dates are not used;
  \item the six-month accounting lag is a conservative approximation rather
        than a report-specific publication date;
  \item annual rather than quarterly Compustat data limit the available
        accounting information;
  \item analyst forecasts and detailed CRSP microstructure variables are
        unavailable;
  \item some predictors are transparent Yahoo/annual-Compustat proxies rather
        than exact replications of every published characteristic;
  \item the raw return target contains large observations that can strongly
        affect squared-error metrics.
\end{itemize}

These limitations do not invalidate the project, but they define the scope of
the empirical conclusions.

The dataset should therefore be interpreted as a transparent and
leakage-controlled course-project panel rather than a complete point-in-time
representation of the entire United States equity market.

% ============================================================
% 20. References
% ============================================================

\section{References}

\begin{thebibliography}{20}

\bibitem{amihud2002}
Amihud, Y. 2002.
Illiquidity and stock returns: Cross-section and time-series effects.
\textit{Journal of Financial Markets}.

\bibitem{ang2006}
Ang, A., Hodrick, R., Xing, Y., and Zhang, X. 2006.
The cross-section of volatility and expected returns.
\textit{Journal of Finance}.

\bibitem{bali2011}
Bali, T., Cakici, N., and Whitelaw, R. 2011.
Maxing out: Stocks as lotteries and the cross-section of expected returns.
\textit{Journal of Financial Economics}.

\bibitem{cboe2019}
Cboe. 2019.
\textit{VIX White Paper}.

\bibitem{fama1992}
Fama, E., and French, K. 1992.
The cross-section of expected stock returns.
\textit{Journal of Finance}.

\bibitem{george2004}
George, T., and Hwang, C. 2004.
The 52-week high and momentum investing.
\textit{Journal of Finance}.

\bibitem{gkx2020}
Gu, S., Kelly, B., and Xiu, D. 2020.
Empirical asset pricing via machine learning.
\textit{Review of Financial Studies}.

\bibitem{green2017}
Green, J., Hand, J., and Zhang, X. 2017.
The characteristics that provide independent information about average
U.S. monthly stock returns.
\textit{Review of Financial Studies}.

\bibitem{jegadeesh1993}
Jegadeesh, N., and Titman, S. 1993.
Returns to buying winners and selling losers.
\textit{Journal of Finance}.

\bibitem{novymarx2013}
Novy-Marx, R. 2013.
The other side of value: The gross profitability premium.
\textit{Journal of Financial Economics}.

\bibitem{sloan1996}
Sloan, R. 1996.
Do stock prices fully reflect information in accruals and cash flows?
\textit{Accounting Review}.

\bibitem{titman2004}
Titman, S., Wei, K., and Xie, F. 2004.
Capital investments and stock returns.
\textit{Journal of Financial and Quantitative Analysis}.

\bibitem{welchgoyal2008}
Welch, I., and Goyal, A. 2008.
A comprehensive look at the empirical performance of equity premium prediction.
\textit{Review of Financial Studies}.

\end{thebibliography}

\end{document}